% ======================================================================
% Chapter 13: Critical Exponents and Scalar Field Theory
% (book pages 439-468; PDF pages 462-491)
% ======================================================================
\chapter{Critical Exponents and Scalar Field Theory}

The idea of running coupling constants and renormalization-group flows gives
us a new language with which to discuss the qualitative behavior of scalar
field theory.  In our first discussion of $\phi^4$ theory, each value of the
coupling constant---and, more generally, each form of the potential and each
spacetime dimension---gave a separate problem to be explored.  But in
Chapter~12, we saw that $\phi^4$ theories with different values of the
coupling are connected by renormalization-group flows, and that the pattern of
these flows changes continuously with the spacetime dimension.  In this
context, it makes sense to ask the very general question: How does $\phi^4$
theory behave as a function of the dimension?  This chapter will give a
detailed answer to this question.

The central ingredient in our analysis will be the Wilson--Fisher fixed point
discussed in Section~12.5.  This fixed point exists in spacetime dimensions
$d$ with $d < 4$; in those dimensions it controls the renormalization group
flows of massless $\phi^4$ theory.  The scalar field theory has manifest or
spontaneously broken symmetry according to the sign of the mass parameter
$m^2$.  Near $m^2 = 0$, the theory exhibits scaling behavior with anomalous
dimensions whose values are determined by the renormalization group
equations.  For $d > 4$, the Wilson--Fisher fixed point disappears, and only
the free-field fixed point remains.  Again, the theory exhibits two distinct
phases, but now the behavior at the transition is determined by the
renormalization group flows near the free-field fixed point, so the scaling
laws are those that follow from simple dimensional analysis.

The continuation of these results to Euclidean space has important
implications for the theory of phase transitions in magnets and fluids.  As we
discussed in the previous chapter, the ideas of the renormalization group
imply that the power-law behaviors of thermodynamic quantities near a phase
transition point are determined by the behavior of correlation functions in a
Euclidean $\phi^4$ theory.  The results stated in the previous paragraph then
imply the following conclusions for critical scaling laws: For statistical
systems in a space of dimension $d > 4$, the scaling laws are just those
following from simple dimensional analysis.  These predictions are precisely
those of Landau theory, which we discussed in Chapter~8.  On the other hand,
for $d < 4$, the critical scaling laws are modified, in a way that we can
compute using the renormalization group.

In $d = 4$, we are on the boundary between the two types of scaling behavior.
This corresponds to the situation in which $\phi^4$ theory is precisely
renormalizable.  In this case, the dimensional analysis predictions are
corrected, but only by logarithms.  We will analyze this case specifically in
Section~13.2.

Though it is not obvious, the case $d = 2$ provides another boundary.  Here
the transition to spontaneous symmetry breaking is described by a different
quantum field theory, which becomes renormalizable in two dimensions.  In
Section~13.3, we will introduce that theory, called the \emph{nonlinear sigma
model}, and show how its renormalization group behavior merges smoothly with
that of $\phi^4$ theory.  By combining all of the results of this chapter, we
will obtain a quantitative understanding of the behavior of $\phi^4$ theory,
and of critical phenomena, over the whole range of spacetime dimensions.

\section{Theory of Critical Exponents}

At the end of Chapter~12, we used properties of the renormalization group for
scalar field theory to make a prediction about the behavior of correlations
near the critical point of a thermodynamic system.  We argued that the range
of correlations, the correlation length $\xi$, should increase to infinity as
one approaches the critical point, according to the scaling law
\eqref{eq:12.141}.  The exponent in this equation, called $\nu$, should
depend only on the symmetry of the order parameter.  We argued, further, that
this exponent is related to the anomalous dimension of a local operator in
$\phi^4$ theory, and that it can be computed from Feynman diagrams.  In this
section, we will show that similar conclusions apply more generally to a
large number of scaling laws associated with a critical point.

To begin, we will define systematically a set of \emph{critical exponents},
exponents of scaling laws that describe the thermodynamic behavior in the
vicinity of the critical point.  We will then show, using the Callan-Symanzik
equation, that all these exponents can be reduced to two basic anomalous
dimensions.  Finally, we will compare this remarkable prediction of quantum
field theory to experiment.

In suggesting a set of critical scaling laws, we begin with the behavior of the
correlation function of fluctuations of the ordering field.  For definiteness,
we will use the language appropriate to a magnet, as in Chapter~8.  We will
compute classical thermal expectation values as correlation functions in a
Euclidean quantum field theory, as explained in Section~9.3.  The fluctuating
field will be called the spin field $s(\mathbf{x})$, its integral will be the
magnetization $M$, and the external field that couples to $s(\mathbf{x})$ will
be called the magnetic field $H$.  (In deference to the magnetization, we will
denote the renormalization scale in the Callan-Symanzik equation by $\mu$ in
this section.)

Define the two-point correlation function by
\begin{equation}
  G(x) = \langle s(x)s(0)\rangle ,
  \label{eq:13.1}
\end{equation}
or by the connected expectation value, if we are in the magnetized phase where
$\langle s(x)\rangle \neq 0$.  Away from the critical point, $G(x)$ should
decay exponentially, according to
\begin{equation}
  G(x) \sim \exp[-|x|/\xi] .
  \label{eq:13.2}
\end{equation}
The approach to the critical point is characterized by the parameter
\begin{equation}
  t = \frac{T - T_c}{T_c} .
  \label{eq:13.3}
\end{equation}
Then we expect that, as $t\to 0$, the correlation length should increase to
infinity.  Define the exponent $\nu$, \eqref{eq:12.141}, by the formula
\begin{equation}
  \xi \sim |t|^{-\nu} .
  \label{eq:13.4}
\end{equation}
Just at $t = 0$, the correlation function should decay only as a power law.
Define the exponent $\eta$ by the formula
\begin{equation}
  G(x) \sim \frac{1}{|x|^{d-2+\eta}} ,
  \label{eq:13.5}
\end{equation}
where $d$ is the Euclidean space dimension.

The behaviors of thermodynamic quantities near the critical point define a
number of additional exponents.  Typically, the specific heat of the
thermodynamic system diverges as $t\to 0$; define the exponent $\alpha$ by the
formula for the specific heat at fixed external field $H = 0$:
\begin{equation}
  C_H \sim |t|^{-\alpha} .
  \label{eq:13.6}
\end{equation}
Since the ordering sets in at $t = 0$, the magnetization at zero field tends
to zero as $t\to 0$ from below.  Define the exponent $\beta$ (not to be
confused with the Callan-Symanzik function) by
\begin{equation}
  M \sim |t|^{\beta} .
  \label{eq:13.7}
\end{equation}
Even at $t = 0$ one has a nonzero magnetization at nonzero magnetic field.
Write the law by which this magnetization tends to zero as $H\to 0$ as the
relation
\begin{equation}
  M \sim H^{1/\delta} .
  \label{eq:13.8}
\end{equation}
Finally, the magnetic susceptibility diverges at the critical point; we write
this divergence as the relation
\begin{equation}
  \chi \sim |t|^{-\gamma} .
  \label{eq:13.9}
\end{equation}
Equations \eqref{eq:13.4}--\eqref{eq:13.9} define a set of critical exponents
$\alpha$, $\beta$, $\gamma$, $\delta$, $\eta$, $\nu$, which can be measured
experimentally for a variety of thermodynamic systems.%
\footnote{A variety of further critical exponents and relations are presented
in M.~E.~Fisher, \emph{Repts.\ Prog.\ Phys.} \textbf{30}, 615 (1967).}

In Chapter~12 we argued, following Wilson, that a thermodynamic system near
its critical point can be described by a Euclidean quantum field theory.  At
the level of the atomic scale, the Lagrangian of this quantum field theory may
be complicated; however, when we have integrated out the small-scale degrees
of freedom, this Lagrangian simplifies.  If we adjust a parameter of the
theory to insure the presence of long-range correlations, the Lagrangian must
closely approach a fixed point of the renormalization group.  Generically, the
Lagrangian will approach the fixed point with a single unstable or relevant
direction, corresponding to the mass parameter of $\phi^4$ theory.  In
$d < 4$, this is the Wilson--Fisher fixed point.  In $d > 4$, it is the
free-field fixed point.  For definiteness, we will assume $d < 4$ in the
following discussion.

\subsection{Exponents of the Spin Correlation Function}

In this setting, we can study the behavior of the spin-spin correlation
function $G(x)$.  By the argument just reviewed, $G(x)$ is proportional to
the two-point correlation function of a Euclidean scalar field theory.  The
technology introduced in the previous chapter can be applied directly.  The
correlation function obeys the Callan-Symanzik equation Eq.~(12.125),
\begin{equation}
  \Bigl[ \mu\frac{\partial}{\partial\mu} + \sum_i \beta_i\frac{\partial}
    {\partial p_i} + \gamma \Bigr]\ G(x; \mu, \{p_i\}) = 0 .
  \label{eq:13.10}
\end{equation}
Here we include the $\phi^4$ coupling $\lambda$ among the generalized
couplings $p_i$.

By dimensional analysis, in $d$ dimensions,
\begin{equation}
  G(x) = \frac{1}{|x|^{d-2}}\ g\bigl( |x|\mu, \{p_i\} \bigr) ,
  \label{eq:13.11}
\end{equation}
where $g$ is an arbitrary function of the dimensionless parameters.  (This is
the Fourier transform of the statement that $G(p) \sim p^{-2}$ times a
dimensionless function.)  From this starting point, we can solve the
Callan-Symanzik equation \eqref{eq:13.10} by the method of Section~12.3, and
find
\begin{equation}
  G(x) = \frac{1}{|x|^{d-2}}\ h\bigl( \{p_i(x)\} \bigr)\ \exp\Bigl[
    \int_{1/|x|}^{\mu} d\log(|x'|)\ \gamma\bigl( \{p(x')\} \bigr) \Bigr] ,
  \label{eq:13.12}
\end{equation}
where $h$ is a dimensionless initial condition.  The running coupling
constants $p_i$ obey the differential equation
\begin{equation}
  \frac{d}{d\log(1/\mu|x|)}\ p_i = \beta_i\bigl( \{p_j\} \bigr) .
  \label{eq:13.13}
\end{equation}
We studied the solution to this equation in Section~12.5.  We saw there that,
for flows that come to the vicinity of the Wilson--Fisher fixed point, the
dimensionless coefficient of the mass operator grows as one moves toward
large distances, while the other dimensionless parameters become small.  Let
$\lambda^{*}$ be the location of the fixed point.  Then we can write more
explicitly
\begin{equation}
  p_m = p_m(\mu|x|)^{-2+\gamma_{\phi^2}(\lambda^{*})},
  \qquad p_i = p_i(\mu|x|)^{-\Delta_i} ,
  \label{eq:13.14}
\end{equation}
where $\Delta_i > 0$ for $i \neq m$.  If the deviation of $\lambda$ from the
fixed point is treated as one of the $p_i$, by defining
\begin{equation}
  p_\lambda = \lambda - \lambda^{*} ,
  \label{eq:13.15}
\end{equation}
this parameter also decreases in importance as a power of $|x|$, as we
demonstrated in Eq.~(12.96).  In the language of Section~12.1, all
of the parameters $p_i$ multiply irrelevant operators, except for $p_m$, which
multiplies a relevant operator.

To approach the critical point, we adjust the parameters of the underlying
theory so that, at some scale $(1/\mu)$ near the atomic scale, $p_m \ll 1$.
If $p_m$ is adjusted by tuning the temperature of the thermodynamic system,
then $p_m \sim t$.  The critical scaling laws will be valid if there is a
region of distance scales where $p_m$ remains small while the other $p_i$ can
be neglected.  The scaling laws can then be computed by evaluating the
solution to the Callan-Symanzik equation with $p_m$ given by \eqref{eq:13.14}
and the other $p_i$ set equal to zero.  The corrections to this approximation
can be shown to be proportional to positive powers of $t$.

In this approximation, we should evaluate the function $\gamma(\lambda)$ in
\eqref{eq:13.12} at $p_\lambda = 0$, that is, at the fixed point.  Using this
value and the solution for $p_m$, Eq.~\eqref{eq:13.12} becomes
\begin{equation}
  G(x) = \frac{1}{|x|^{d-2+2\gamma(\lambda^{*})}}\ h\bigl(
    p_m|x|^{2-\gamma_{\phi^2}(\lambda^{*})} \bigr) .
  \label{eq:13.16}
\end{equation}
This equation implies the scaling laws \eqref{eq:13.5} and \eqref{eq:13.4}:
For the argument of $h$ sufficiently small, $G(x)$ obeys Eq.~\eqref{eq:13.5},
with
\begin{equation}
  \eta = 2\gamma(\lambda^{*}) .
  \label{eq:13.17}
\end{equation}
At large distances, $h$ must fall off exponentially, since this function is
derived from a scalar field propagator.  From the argument of $h$, we deduce
that this exponential must be of the form
\begin{equation}
  \exp\Bigl[ -|x|\bigl( p_m \bigr)^{\nu} \Bigr] ,
  \label{eq:13.18}
\end{equation}
where, as in \eqref{eq:12.139},
\begin{equation}
  \nu^{-1} = 2 - \gamma_{\phi^2}(\lambda^{*}) .
  \label{eq:13.19}
\end{equation}
This is precisely the scaling law \eqref{eq:13.2}, \eqref{eq:13.4}, with the
identification of $\nu$ in terms of the anomalous dimension of the operator
$\phi^2$.

\subsection{Exponents of Thermodynamic Functions}

The thermodynamic critical scaling laws can be derived in a similar way, by
studying the scaling behavior of macroscopic thermodynamic variables.  These
are derived from the Gibbs free energy, or, in the language of quantum fields,
from the effective potential of the scalar field theory.  Since the effective
potential, and, more generally, the effective action, are constructed from
correlation functions, the scaling behavior of these quantities can be derived
from the scaling laws for correlation functions.

The Gibbs free energy $G(M, t)$ is related to the effective potential
$V_{eff}(\phi_{cl})$ by identifying the classical field $\phi_{cl}$ with the
magnetization $M$ and the external field $H$ with the derivative of $G$ with
respect to $M$.  The scaling behavior of $G$ follows from the fact that it is
the generating functional of correlation functions.  In the scaling region,
one obtains the scaling formula
\begin{equation}
  G(M, t) = M^{1+\delta}\ h\bigl( t\,M^{-1/\beta} \bigr) ,
  \label{eq:13.35}
\end{equation}
where $h$ has a smooth limit as $t\to 0$.  An equivalent way to represent this
formula is
\begin{equation}
  G(M, t) = t^{\beta(1+\delta)}\ f\bigl( M\,t^{-\beta} \bigr) .
  \label{eq:13.36}
\end{equation}
The scaling laws for thermodynamic quantities follow immediately from these
relations.  Along the line $t = 0$, we find from \eqref{eq:13.35} that
\begin{equation}
  H = \frac{\partial G}{\partial M} = h(0)\,M^{\delta} ,
  \label{eq:13.37}
\end{equation}
which is precisely \eqref{eq:13.8}.  Below the critical temperature, we find
the nonzero value of the magnetization by minimizing $G$ with respect to
$M$.  In the scaling region, this minimum occurs at the minimum $m_0$ of the
function $f(m)$ in \eqref{eq:13.36}.  This leads to relation \eqref{eq:13.7},
in the form
\begin{equation}
  M\,t^{-\beta} = m_0 .
  \label{eq:13.38}
\end{equation}
If we work above $T_c$ and in zero field, the minimum of $f$ must occur at
$M = 0$.  Then
\begin{equation}
  G(t) \sim t^{\beta(1+\delta)} .
  \label{eq:13.39}
\end{equation}
To compute the specific heat, we differentiate twice with respect to
temperature; this gives the scaling law \eqref{eq:13.6}, with
\begin{equation}
  2 - \alpha = \beta(1+\delta) = \frac{d}{2 - \gamma_{\phi^2}(\lambda^{*})} .
  \label{eq:13.40}
\end{equation}
Finally, we must construct the scaling law for the magnetic susceptibility.
From \eqref{eq:13.36}, the scaling law for $H$ at nonzero $t$ is
\begin{equation}
  H = \frac{\partial G}{\partial M} = t^{\beta\delta}\ f'\bigl( M\,t^{-\beta}
    \bigr) .
  \label{eq:13.41}
\end{equation}
The inverse of this relation is the scaling law
\begin{equation}
  M = t^{\beta}\ c\bigl( H\,t^{-\beta\delta} \bigr) .
  \label{eq:13.42}
\end{equation}
The magnetic susceptibility at zero field is then
\begin{equation}
  \chi = \Bigl( \frac{\partial M}{\partial H} \Bigr)_t
  = c'(0)\,t^{-\beta(\delta-1)} .
  \label{eq:13.43}
\end{equation}
Thus, we confirm Eq.~\eqref{eq:13.9}, with the identification
\begin{equation}
  \gamma = (\delta-1)\beta = 2\bigl( 1 - \gamma_{\phi^2}(\lambda^{*}) \bigr) .
  \label{eq:13.44}
\end{equation}

We have now found explicit expressions for all of the various critical
exponents in terms of the Callan-Symanzik functions.  As the dimensionality
$d$ approaches 4 from below, the fixed point $\lambda^{*}$ tends to zero.
Then the six critical exponents approach the values that they would attain in
simple dimensional analysis:
\begin{equation}
  \eta = 0;\qquad \nu = \tfrac{1}{2};\qquad \alpha = 0;\qquad
  \beta = \tfrac{1}{2};\qquad \gamma = 1;\qquad \delta = 3 .
  \label{eq:13.45}
\end{equation}
It is no surprise that the values of $\eta$, $\nu$ and $\beta$ given in
\eqref{eq:13.45} are those that we derived in Chapter~8 from the Landau
theory of critical phenomena.  The other values can similarly be shown to
follow from Landau theory.  The renormalization group analysis tells us how
to systematically correct the predictions of Landau theory to take proper
account of the large-scale fluctuations of the spin field.

Notice that all of the exponents associated with thermodynamic quantities are
constructed from the same ingredients as the exponents associated with the
correlation function.  From the field theory viewpoint, this is obvious,
since all of the scaling laws in the field theory must ultimately follow from
the anomalous dimensions of the operators $\phi^2(x)$ and $\phi^4(x)$, which
are precisely $\gamma_{\phi^2}(\lambda^{*})$ and $\gamma(\lambda^{*})$.  This
result, however, has an interesting experimental consequence: It implies
\emph{model-independent relations} among critical exponents.  For example, in
any system with a critical point, this theory predicts
\begin{equation}
  \alpha = 2 - d\nu,\qquad \beta = \tfrac{1}{2}(d-2+\eta)\nu .
  \label{eq:13.46}
\end{equation}
These relations test the general framework of identifying a critical point
with the fixed point of a renormalization group flow.

In addition, the field theoretic approach to critical phenomena predicts that
critical exponents are \emph{universal}, in the sense that they take the same
values in condensed matter systems that approach the same scalar field fixed
point in the limit $T\to T_c$.

\subsection{Values of the Critical Exponents}

Finally, scalar field theory actually predicts the values of
$\gamma(\lambda^{*})$ and $\gamma_{\phi^2}(\lambda^{*})$, either from the
expansion in powers of $\epsilon = 4-d$ described in Section~12.5 or by
direct expansion of the $\beta$ and $\gamma$ functions in powers of
$\lambda$.  We can use these expressions to generate quantitative predictions
for the critical exponents.  We gave an example of such a prediction at the
end of Section~12.5, when we presented in Eq.~\eqref{eq:12.143} the first two
terms of an expansion for $\nu$.  We now return to this question to give
field-theoretic predictions for all of the critical exponents.

In our discussion at the end of Section~12.5, we remarked that magnets with
different numbers of fluctuating spin components are observed to have
different values for the critical exponents.  An optimistic hypothesis would
be that any thermodynamic system with $N$ fluctuating spin components, or,
more generally, $N$ fluctuating thermodynamic variables at the critical
point, would be described by the same fixed point field theory with $N$
scalar fields.  A natural candidate for this fixed point would be the
Wilson--Fisher fixed point of the $O(N)$-symmetric $\phi^4$ theory discussed
in Chapter~11.  We will now describe the computation of critical exponents
for general values of $N$ in this theory.

As a first step, we should compute the values of the functions
$\beta(\lambda)$, $\gamma(\lambda)$, and $\gamma_{\phi^2}(\lambda)$ in four
dimensions.  This computation parallels the analysis done in Chapter~12 for
ordinary $\phi^4$ theory, so we will only indicate the changes that need to
be made for this case.  Just as in ordinary $\phi^4$ theory, the propagator of
the massless $O(N)$-symmetric theory receives no field strength corrections
in one-loop order, and so the one-loop term in $\gamma(\lambda)$ again
vanishes.  In Problem~13.2, we compute the leading, two-loop, contribution to
$\gamma(\lambda)$ in $O(N)$-symmetric $\phi^4$ theory:
\begin{equation}
  \gamma = (N+2)\ \frac{\lambda^2}{(4\pi)^4}\ \frac{1}{6}
    + \mathcal{O}(\lambda^3) .
  \label{eq:13.47}
\end{equation}
The one-loop contribution to the $\beta$ function in $\phi^4$ theory is
derived from the one-loop vertex counterterm $\delta_\lambda$, given in
Eq.~(12.44).  For the $O(N)$-symmetric case, we computed the
divergent part of the corresponding vertex counterterm in Section~11.2; from
Eq.~\eqref{eq:11.22},
\begin{equation}
  \delta_\lambda = \lambda^2(N+8)\ \frac{\Gamma(2-d/2)}{(4\pi)^2}
    + \cdots .
  \label{eq:13.48}
\end{equation}
Following the logic to Eq.~\eqref{eq:12.46}, or using Eq.~\eqref{eq:12.54},
we find
\begin{equation}
  \beta = (N+8)\ \frac{\lambda^2}{16\pi^2} + \mathcal{O}(\lambda^3) .
  \label{eq:13.49}
\end{equation}
This reduces to the $\beta$ function of $\phi^4$ theory if we set $N = 1$ and
replace $\lambda\to\lambda/6$, as indicated below Eq.~\eqref{eq:11.5}.
Finally, to compute $\gamma_{\phi^2}$, we must repeat the computation done at
the end of Section~12.4.  If we consider, instead of \eqref{eq:12.113}, the
Green's function $\langle\phi^i(p)\phi^j(q)\phi^2(k)\rangle$, and replace the
vertex of $\phi^4$ theory by the four-point vertex following from the
Lagrangian \eqref{eq:11.5}, the factor $(-i\lambda)$ in the first line of
\eqref{eq:12.114} is replaced by
\begin{equation}
  (-2i\lambda)\bigl( \delta^{ij}\delta^{kl} + \delta^{ik}\delta^{jl}
    + \delta^{il}\delta^{jk} \bigr)\cdot\delta^{kl}
  = -2i\lambda(N+2)\delta^{ij} .
\end{equation}
Then
\begin{equation}
  \gamma_{\phi^2} = (N+2)\ \frac{\lambda}{16\pi^2} + \mathcal{O}(\lambda^2) .
  \label{eq:13.50}
\end{equation}

Next, we consider the same theory in $(4-\epsilon)$ dimensions.  The $\beta$
function now becomes
\begin{equation}
  \beta = -\epsilon\lambda + (N+8)\ \frac{\lambda^2}{16\pi^2} ,
  \label{eq:13.51}
\end{equation}
so there is a Wilson--Fisher fixed point at
\begin{equation}
  \lambda^{*} = \frac{16\pi^2\epsilon}{N+8} .
  \label{eq:13.52}
\end{equation}
At this fixed point, we find
\begin{equation}
  \gamma(\lambda^{*}) = \frac{N+2}{4(N+8)^2}\,\epsilon^2 + \cdots,\qquad
  \gamma_{\phi^2}(\lambda^{*}) = \frac{N+2}{N+8}\,\epsilon + \cdots .
  \label{eq:13.53}
\end{equation}
From these two results, we can work out predictions for the whole set of
critical exponents to order $\epsilon$.  As an example, inserting
\eqref{eq:13.53} into \eqref{eq:13.19}, we find
\begin{equation}
  \nu^{-1} = 2 - \frac{N+2}{N+8}\,\epsilon + \mathcal{O}(\epsilon^2) ,
  \label{eq:13.54}
\end{equation}
as claimed at the end of Section~12.5.

In our discussion in Chapter~12, we claimed that the predictions of critical
exponents are in rough agreement with experimental data.  However, by
computing to higher order, one can obtain a much more precise comparison of
theory and experiment.  The $\epsilon$ expansion of critical exponents has now
been worked out through order $\epsilon^5$.  More impressively, the
$\lambda$ expansion for critical exponents in $d = 3$ has been worked out
through order $\lambda^9$.  By summing this perturbation series, it is
possible to obtain very precise estimates of the anomalous dimensions
$\gamma(\lambda^{*})$ and $\gamma_{\phi^2}(\lambda^{*})$ and, through them,
precise predictions for the critical exponents.

A comparison of these values to direct determinations of the critical
exponents is given in Table~\ref{tab:13.1}.  The column labeled ``QFT'' gives
values of critical exponents obtained by anomalous dimension calculations
using $\phi^4$ perturbation theory in three dimensions.  The column labeled
``Experiment'' lists a selection of experimental determinations of the
critical exponents in a variety of systems.  These include the liquid-gas
critical point in Xe, CO$_2$, and other fluids, the critical point in binary
fluid mixtures with liquid-liquid phase separation, the order-disorder
transition in the atomic arrangement of the Cu-Zn alloy $\beta$-brass, the
superfluid transition in $^4$He, and the order-disorder transitions in
ferromagnets (EuO, EuS, Ni) and antiferromagnets (RbMnF$_3$).  The agreement
between experimental determinations of the exponents in different systems is
a direct test of universality.  For the case of systems with a single order
parameter ($N = 1$), there is a remarkable diversity of physical systems that
are characterized by the same critical exponents.

\begin{table}[t]
  \centering
  \caption{Critical exponents in three dimensions: comparison of QFT
    predictions, lattice estimates, and experiment.}
  \begin{tabular}{cccc}
    \toprule
    & QFT & Lattice & Experiment\\
    \midrule
    $N=1$: $\nu$ & 0.630 & 0.630 & 0.63\\
    \quad $\eta$ & 0.033 & 0.031 & 0.03\\
    $N=2$: $\nu$ & 0.670 & 0.670 & 0.67\\
    \quad $\eta$ & 0.038 & 0.035 & 0.04\\
    $N=3$: $\nu$ & 0.704 & 0.704 & 0.71\\
    \quad $\eta$ & 0.041 & 0.040 & 0.04\\
    \bottomrule
  \end{tabular}
  \label{tab:13.1}
\end{table}

The column labeled ``Lattice'' contains estimates of critical exponents in
abstract lattice statistical mechanical models.  For these simplified models,
the statistical mechanical partition function can be calculated in an
expansion for large temperature.  With some effort, these expansions can be
carried out to high order and resummed, giving very precise values for the
critical exponents.  The agreement among the three columns of
Table~\ref{tab:13.1}---field theory, lattice, and experiment---is a striking
confirmation of the renormalization group description of critical phenomena.

\section{Critical Behavior in Four Dimensions}

In this section we study the case $d = 4$, where $\phi^4$ theory is precisely
renormalizable.  Here the running coupling constant of massless $\phi^4$
theory tends to zero logarithmically at large distances, but grows at large
momenta.  The dimensional-analysis predictions of Landau theory are corrected
by powers of logarithms.  The most interesting application of this analysis
is the resolution of the problem posed at the end of Section~11.4: the
question of whether the effective potential of $\phi^4$ theory at $\mu^2 = 0$
has a symmetry-breaking minimum.

We begin with the one-loop expression for the effective potential in
$\phi^4$ theory with $\mu^2 = 0$.  From Eq.~\eqref{eq:11.81}, the
$\overline{MS}$ expression is
\begin{equation}
  V_{eff}(\phi_{cl}) = \frac{\lambda}{4}\phi_{cl}^4\ \Bigl[ 1 +
    \frac{\lambda}{16\pi^2}\ \Bigl( 9\Bigl( \log\frac{\phi_{cl}^2}{M^2}
      - \frac{25}{6} \Bigr) \Bigr) \Bigr] ,
  \label{eq:13.55}
\end{equation}
where we have set $N = 1$.  This expression has an apparent minimum at a
nonzero value of $\phi_{cl}$, but the minimum is an artifact of the
perturbation expansion: it occurs where the logarithm is large enough to
cancel the leading term, so that perturbation theory breaks down.

To resolve this question, we apply the renormalization group.  The running
coupling constant $\bar{\lambda}(p)$ satisfies
\begin{equation}
  \frac{d}{d\log(p/M)}\ \bar{\lambda} = \frac{3\bar{\lambda}^2}{16\pi^2},
  \qquad \bar{\lambda}(M) = \lambda .
\end{equation}
The improved effective potential is obtained by evaluating the one-loop
formula with the running coupling and field evaluated at the scale set by
$\phi_{cl}$:
\begin{equation}
  V_{eff}(\phi_{cl}) = \frac{\bar{\lambda}(\phi_{cl})}{4}\ \bar{\phi}_{cl}^4
    \ \cdots ,
  \label{eq:13.64}
\end{equation}
where the running field $\bar{\phi}_{cl}$ incorporates the anomalous
dimension.  The error in Eq.~\eqref{eq:13.64} comes in the determination of
$\bar{\lambda}$ as a power series in $\lambda$.  Thus this error is of order
$\lambda^3$.  As $\phi_{cl}\to 0$, $\lambda\to 0$, and so the representation
\eqref{eq:13.64} becomes more and more accurate.  Thus this formula
successfully sums the powers of the dangerous logarithm.  Viewed as a
function of $\phi_{cl}$, \eqref{eq:13.64} has its minimum at
$\phi_{cl} = 0$.  Thus the apparent symmetry-breaking minimum of
\eqref{eq:13.55} is indeed an artifact of the incomplete perturbation
expansion and disappears in a more complete treatment.  This resolves the
question that we raised in Section~11.4.  We should note that, in more
complicated examples, an apparent symmetry-breaking minimum of the effective
potential found in the one-loop order of perturbation theory can survive a
renormalization-group analysis.  An example is given in the Final Project for
Part~II.

The procedure we have followed in this argument is called the
\emph{renormalization group improvement} of perturbation theory.  The
technique can be applied equally well to the computation of correlation
functions and other predictions of Feynman diagram perturbation theory: One
compares the solution of the Callan-Symanzik equation to the result of a
straightforward perturbation theory computation to the same order in the
coupling constant, choosing the undetermined function in the renormalization
group solution in such a way as to reproduce the perturbation theory result.
In this way, one finds a more compact formula in which large logarithms such
as those in Eq.~(13.56) are resummed into running coupling constants.
This resummation produces the dependence of correlation functions on the
logarithm of the mass scale that characterizes a field theory with a marginal
or renormalizable perturbation.

In the case of $\phi^4$ theory, the running coupling constant goes to zero at
small momenta and becomes large at large momenta.  Since the error term in
improved perturbation theory is a power of $\lambda$, the improved
perturbation theory becomes accurate at small momenta but goes out of control
at large momenta.  This accords with our physical intuition: We would expect
perturbation theory to be accurate only when the running coupling constant
stays small.

In an asymptotically free theory, where the running coupling constant becomes
small at large momenta, we can find accurate expressions for correlation
functions at large momenta using renormalization-group-improved perturbation
theory.  In Chapters~17 and~18 we will use this idea as our major tool in
analyzing the short-distance behavior of the strong interactions.

\section{The Nonlinear Sigma Model}

To complete our study of scalar field theory, we will discuss a nonlinear
theory of scalar fields, whose structure is very different from that of
$\phi^4$ theory.  This theory, called the \emph{nonlinear sigma model}, was
first proposed as an alternative description of spontaneous symmetry
breaking.  It will be interesting to us for three reasons.  First, it provides
a simple explicit example of an asymptotically free theory.  Second, it will
give us a second dimensional expansion with which we can study the
Wilson--Fisher fixed point.  Then we can see where the Wilson--Fisher fixed
point goes in the space of Lagrangians for dimensions $d$ well below 4.
Finally, we will show that the nonlinear sigma model is exactly solvable in a
limit that is different from the standard weak-coupling limit.  This solution
will give us further insight into the dependence of symmetry breaking on
spacetime dimensionality.

\subsection{The $d = 2$ Nonlinear Sigma Model}

We begin our study in two dimensions.  In $d = 2$, a scalar field is
dimensionless; thus, any theory of scalar fields $\phi^i$ with a Lagrangian
of the form
\begin{equation}
  \mathcal{L} = \tfrac{1}{2}\sum_i f\bigl( \{\phi^j\} \bigr)\,
    (\partial_\mu\phi^i)^2 ,
  \label{eq:13.65}
\end{equation}
has dimensionless couplings and so is renormalizable.  Since any function
$f$ leads to a renormalizable theory, this class of scalar field theories
contains an infinite number of marginal parameters.  To restrict these
possible parameters, we must impose some symmetries on the theory.

A simple choice is to take the scalar fields $\phi^i$ to form an
$N$-component unit vector field $n^i(x)$, constrained to satisfy
\begin{equation}
  \sum_{i=1}^{N} |n^i(x)|^2 = 1 .
  \label{eq:13.66}
\end{equation}
If we insist that the field theory has $O(N)$ symmetry, the function $f$ in
\eqref{eq:13.65} can depend only on the invariant length of $n(x)$, which is
constrained by \eqref{eq:13.66}.  Thus, the most general possible choice for
$f$ is a constant.  Similarly, the only possible nonderivative interaction
$g(\{n^i\})$ that one might add to \eqref{eq:13.65} is a constant, and this
would have no effect on the Green's functions of $n$.  With these
restrictions, the most general Lagrangian one can build from $n(x)$ with two
derivatives and $O(N)$ symmetry is
\begin{equation}
  \mathcal{L} = \frac{1}{2g_0^2}\ (\partial_\mu n)^2 .
  \label{eq:13.67}
\end{equation}
This theory has a straightforward physical interpretation.  It is a
phenomenological description of a system with $O(N)$ symmetry spontaneously
broken by the vacuum expectation value of a field that transforms as a vector
of $O(N)$.  Consider, for example, the situation in $N$-component $\phi^4$
theory in its spontaneously broken phase.  The field $\phi^i$ acquires a
vacuum expectation value, which we can write in terms of a magnitude and a
direction parameterized by a unit vector
\begin{equation}
  \langle\phi^i\rangle = \phi_0\,n^i(x) .
  \label{eq:13.68}
\end{equation}
The fluctuations of $\phi_0$ correspond to a massive field, the field called
$\sigma$ in Chapter~11.  The fluctuations of the direction of the unit vector
$n(x)$ correspond to the $N-1$ Goldstone bosons.  Notice that $n$ has $N$
components subject to the one constraint \eqref{eq:13.66}, and so contains
$N-1$ degrees of freedom.  Formally, the nonlinear sigma model is the limit
of $\phi^4$ theory as the mass of the $\sigma$ field is taken to infinity
while $\phi_0$ is held constant.

Despite this suggestive connection, we will first analyze the nonlinear sigma
model on its own footing as an independent quantum field theory.  It is
convenient to solve the constraint and parametrize $n$ by $N-1$ Goldstone
boson fields $\pi^k$:
\begin{equation}
  n^i = (\pi^1, \ldots, \pi^{N-1}, \sigma) ,
  \label{eq:13.69}
\end{equation}
where, by definition,
\begin{equation}
  \sigma = (1-\pi^2)^{1/2} .
  \label{eq:13.70}
\end{equation}
The configuration $\pi^k = 0$ corresponds to a uniform state of spontaneous
symmetry breaking, oriented in the $N$ direction.  The representation
\eqref{eq:13.69} implies that
\begin{equation}
  (\partial_\mu n)^2 = (\partial_\mu\pi)^2 + \frac{(\pi\cdot\partial_\mu\pi)^2}
    {1-\pi^2} .
  \label{eq:13.71}
\end{equation}
Then the Lagrangian \eqref{eq:13.67} takes the form
\begin{equation}
  \mathcal{L} = \frac{1}{2g_0^2}\ \Bigl[ (\partial_\mu\pi)^2
    + \frac{(\pi\cdot\partial_\mu\pi)^2}{1-\pi^2} \Bigr] .
  \label{eq:13.72}
\end{equation}
Notice that there is no mass term for the field $\pi$, as required by
Goldstone's theorem.

The perturbation theory for the $\pi^k$ field can be read off
straightforwardly by expanding the Lagrangian in powers of $\pi^k$:
\begin{equation}
  \mathcal{L} = \frac{1}{2g_0^2}(\partial_\mu\pi)^2
    + \frac{1}{2g_0^2}(\pi\cdot\partial_\mu\pi)^2 + \cdots .
  \label{eq:13.73}
\end{equation}
This leads to the Feynman rules shown in Fig.~\ref{fig:13.1}, plus additional
vertices with all even numbers of $\pi^k$ fields.  Since the Lagrangian
\eqref{eq:13.67} is the most general $O(N)$-symmetric Lagrangian with
dimensionless coefficients that can be built out of these fields, the theory
can be made finite by renormalization of the coupling constant $g$ and
$O(N)$-symmetric rescaling of the fields $\pi^k$ and $\sigma$.  In
renormalized perturbation theory, there are divergences and counterterms for
each possible $2n$-$\pi$ vertex; however, these counterterms are all related
by the basic requirement that the bare Lagrangian preserve the $O(N)$
symmetry.

\begin{figure}[t]
  \centering
  \includegraphics[width=0.6\textwidth]{fig13.1.png}
  \caption{Feynman rules for the nonlinear sigma model.}
  \label{fig:13.1}
\end{figure}

We now compute the Callan-Symanzik functions for this theory.  Since the
theory is renormalizable, its Green's functions obey the Callan-Symanzik
equation for some functions $\beta$, $\gamma$.  Explicitly,
\begin{equation}
  \Bigl[ M\frac{\partial}{\partial M} + \beta(g)\frac{\partial}{\partial g}
    + n\gamma(g) \Bigr]\ G^{(n)} = 0 ,
  \label{eq:13.74}
\end{equation}
where $G^{(n)}$ is a Green's function of $n$ fields $\pi^k$ or $\sigma$.  To
identify the $\beta$ and $\gamma$ functions, to the leading order in
perturbation theory, we compute two simple Green's functions to one-loop
order and then see what forms are necessary if the Callan-Symanzik equation
is to be satisfied.

The first Green's function we consider is
\begin{equation}
  G^{(1)} = \langle\sigma(x)\rangle .
  \label{eq:13.75}
\end{equation}
Expanding the definition \eqref{eq:13.70}, we find
\begin{equation}
  \langle\sigma(0)\rangle = \Bigl\langle 1 - \frac{1}{2}\pi^2(0) + \cdots
    \Bigr\rangle = 1 - \frac{1}{2}\langle\pi^2(0)\rangle + \cdots .
  \label{eq:13.76}
\end{equation}
To evaluate this formula, we use the propagator of Fig.~\ref{fig:13.1} to
compute
\begin{equation}
  \langle\pi^k(0)\pi^\ell(0)\rangle = (N-1)\,\delta^{k\ell}\ \int
    \frac{d^dk}{(2\pi)^d}\ \frac{g_0^2}{k^2+\mu^2} .
  \label{eq:13.77}
\end{equation}
We have added a small mass $\mu$ as an infrared cutoff.  Then
\begin{equation}
  \langle\pi^k(0)\pi^\ell(0)\rangle = (N-1)\ \frac{g_0^2}{(4\pi)^{d/2}}\
    \Gamma(d/2-1)\,\mu^{2-d}\,\delta^{k\ell} .
  \label{eq:13.78}
\end{equation}
Using this result in our expression for $\langle\sigma\rangle$ and then
subtracting at the momentum scale $M$, we find
\begin{equation}
  \langle\sigma(0)\rangle = 1 - \frac{N-1}{2}\ \frac{g^2}{4\pi}\
    \log\frac{M}{\mu} + \cdots .
  \label{eq:13.79}
\end{equation}
This expression satisfies the Callan-Symanzik equation to order $g$ only if
\begin{equation}
  \gamma(g) = \frac{N-1}{4\pi}\,g^2 + \mathcal{O}(g^4) .
  \label{eq:13.80}
\end{equation}

Next, consider the $\pi^k$ two-point function,
\begin{equation}
  \langle\pi^k(p)\pi^\ell(-p)\rangle = \frac{i\delta^{k\ell}}{p^2}
    + \frac{i\delta^{k\ell}}{p^2}\ \Pi^{k\ell}(p)\ \frac{i}{p^2} + \cdots ,
  \label{eq:13.81}
\end{equation}
where the one-loop self-energy $\Pi^{k\ell}(p)$ is computed from the Feynman
rules in Fig.~\ref{fig:13.1}.  In evaluating $\Pi^{k\ell}$ we again encounter
the integral \eqref{eq:13.77}, and also the integral
\begin{equation}
  \int \frac{d^dk}{(2\pi)^d}\ \frac{i g_0^2\,k^2}{(k^2+\mu^2)^2}
  = \frac{i g_0^2\,d}{(4\pi)^{d/2}}\ \Gamma(2-d/2)\,\mu^{d-2}\ \cdots .
  \label{eq:13.82}
\end{equation}
This formula has no pole at $d = 0$, and for $d > 0$ it is proportional to a
positive power of $\mu^2$; hence, we can set this contraction to zero.  Then
\begin{equation}
  \Pi^{k\ell}(p) = \cdots .
  \label{eq:13.83}
\end{equation}
Subtracting at $M$ as above and taking the limit $d\to 2$, we find
\begin{equation}
  \langle\pi^k(p)\pi^\ell(-p)\rangle = \frac{i\delta^{k\ell}}{p^2}\ \Bigl[
    1 - \frac{N-1}{4\pi}\,g^2\ \log\frac{M^2}{-p^2} \Bigr] .
  \label{eq:13.84}
\end{equation}
Applying the Callan-Symanzik equation to this result gives
\begin{equation}
  \Bigl[ M\frac{\partial}{\partial M} + \beta(g)\frac{\partial}{\partial g}
    + 2\gamma(g) \Bigr]\ \langle\pi^k(p)\pi^\ell(-p)\rangle = 0 ,
\end{equation}
from which we can determine $\beta(g)$:
\begin{equation}
  \beta(g) = -\frac{N-2}{4\pi}\,g^3 + \mathcal{O}(g^5) .
  \label{eq:13.86}
\end{equation}
The sign of the $\beta$ function is negative for $N > 2$.  This is the
signature of \emph{asymptotic freedom}: the running coupling constant
decreases at large momenta.  The nonlinear sigma model in two dimensions is
our first example of an asymptotically free theory.

The negative sign of $\beta$ means that, at large distances, the coupling
constant $g^2$ grows, and the perturbation expansion breaks down.  This is
the statement that the ordering in the nonlinear sigma model is destroyed at
long distances, in accord with the Mermin--Wagner theorem.  To understand the
large-distance behavior more completely, we can integrate the $\beta$
function:
\begin{equation}
  \bar{g}^2(p) = \frac{16\pi^2}{N-2}\ \frac{1}{\log(p/\Lambda)} ,
  \label{eq:13.87}
\end{equation}
where $\Lambda$ is a dynamically generated scale.  The coupling constant
diverges at $p = \Lambda$, signaling the breakdown of perturbation theory and
the disappearance of the ordered phase at long distances.

\subsection{The Nonlinear Sigma Model for $2 < d < 4$}

We now extend the results of this analysis to dimensions $d > 2$.  In general
$d$, we will continue to define the action of the nonlinear sigma model by
\begin{equation}
  \int d^dx\ \mathcal{L} = \int d^dx\ \frac{1}{2g_0^2}\ (\partial_\mu n)^2 ,
  \label{eq:13.102}
\end{equation}
where $n$ is still dimensionless, since it obeys the constraint
$|n|^2 = 1$.  Thus $g$ has the dimensions
$(\mathrm{mass})^{(2-d)/2}$.  We define a dimensionless coupling by writing
\begin{equation}
  T = g^2 M^{d-2} ,
  \label{eq:13.103}
\end{equation}
just as we did in Eq.~(12.122).  If \eqref{eq:13.102} is viewed as the
Boltzmann weight of a partition function, then $T$ is a dimensionless
variable proportional to the temperature.

From \eqref{eq:13.103}, we can find the $\beta$ function for $T$ in $d$
dimensions, in analogy to Eq.~\eqref{eq:12.131}:
\begin{equation}
  \beta(T) = (d-2)T + 2g\beta^{(2)}(g) ,
  \label{eq:13.104}
\end{equation}
where the factor of $2g$ in the second term comes from the definition
$T \sim g^2$.  Since $n$ is dimensionless, the $\gamma$ function is unchanged
from the two-dimensional result when expressed in terms of dimensionless
couplings.  Thus, in $d = 2 + \epsilon$,
\begin{equation}
  \beta(T) = +\epsilon T - (N-2)\ \frac{T^2}{4\pi};\qquad
  \gamma(T) = (N-1)\ \frac{T}{4\pi} .
  \label{eq:13.105}
\end{equation}
Notice that the $\beta$ function for $T$ has a nontrivial zero, which
approaches $T = 0$ as $\epsilon\to 0$.  This zero is located at
\begin{equation}
  T^{*} = \frac{4\pi\epsilon}{N-2} .
  \label{eq:13.106}
\end{equation}
The form of the $\beta$ function is sketched in Fig.~\ref{fig:13.3}.  In
contrast to the Wilson--Fisher zero in $d = 4-\epsilon$, discussed in
Section~12.5, this is an \emph{ultraviolet-stable} fixed point.  The flows to
the infrared go out from this fixed point.  Since $T$ is proportional to the
temperature of the corresponding statistical system, $t\to 0$ is the state of
complete order, while $t\to 1$ is the state of complete disorder.  This
agrees with the intuition that accompanied Polyakov's derivation of the
$\beta$ function.  The fixed point $T^{*}$ corresponds to the critical
temperature.  Thus, the critical temperature tends to zero as $d\to 2$, in
accord with the Mermin--Wagner theorem.

\begin{figure}[t]
  \centering
  \includegraphics[width=0.45\textwidth]{fig13.3.png}
  \caption{The form of the $\beta$ function in the nonlinear sigma model for
    $d > 2$.}
  \label{fig:13.3}
\end{figure}

We can now compute the critical exponents of the nonlinear sigma model in an
expansion in $\epsilon = d-2$.  The exponent $\eta$ is given straightforwardly
by
\begin{equation}
  \eta = \gamma(T^{*}) = (N-1)\ \frac{\epsilon}{N-2} .
  \label{eq:13.107}
\end{equation}
To find the second exponent $\nu$, we need to identify the relevant
perturbation that corresponds to the renormalization group flow away from the
fixed point for $T \neq T_c$.  This is just the deviation of $T$ from
$T^{*}$:
\begin{equation}
  p_T = T - T^{*} .
  \label{eq:13.108}
\end{equation}
From the renormalization group equation for the running coupling constant, we
find that the running $p_T$ obeys
\begin{equation}
  \frac{d}{d\log(p/M)}\ p_T = \Bigl( -\epsilon + 2\frac{N-2}{4\pi}T^{*}
    \Bigr)\,p_T .
  \label{eq:13.109}
\end{equation}
The quantity in brackets is negative.  As in Eqs.~\eqref{eq:12.134} and
\eqref{eq:12.137}, we can identify this quantity with $(-1/\nu)$: At a
momentum $p \ll M$,
\begin{equation}
  p_T(p) \sim p^{1/\nu} ,
  \label{eq:13.110}
\end{equation}
thus $p(p)$ becomes of order 1 at a momentum that is the inverse of
$\xi \sim (T-T^{*})^{-\nu}$, as required.  Using the explicit form of the
$\beta$ function from \eqref{eq:13.105}, we find
\begin{equation}
  \nu^{-1} = \epsilon ,
  \label{eq:13.111}
\end{equation}
independent of $N$ to this order in $\epsilon$.  (Of course, these results
apply only for $N > 2$.)  The thermodynamic critical exponents can be found
from \eqref{eq:13.107} and \eqref{eq:13.111} using the model-independent
relations derived in Section~13.1.  When the values found here for $\nu$ and
$\eta$ are extrapolated to $d = 3$ (that is, $\epsilon = 1$), the agreement
with experiment is not spectacular, but the results at least suggest that the
fixed point we have found here may be the continuation of the Wilson--Fisher
fixed point to the vicinity of two dimensions.

\subsection{Exact Solution for Large $N$}

It is possible to obtain further insight into the nature of this fixed point
by attacking the nonlinear sigma model using another approach.  Since the
nonlinear sigma model depends on a parameter $N$, the number of components of
the unit vector, it is reasonable to ask how this model behaves as
$N\to\infty$.  We now show that if we take this limit holding $g^2N$ fixed,
we can obtain an exact solution to the model with nontrivial behavior.

The manipulations that lead to this solution are most clear if we work in
Euclidean space, regarding the Lagrangian as the Boltzmann weight of a spin
system.  Then we must compute the functional integral
\begin{equation}
  Z = \int \mathcal{D}n\ \exp\Bigl[ -\int d^dx\ \frac{1}{2g_0^2}
    (\partial_\mu n)^2 \Bigr]\ \prod_x \delta\bigl( n^2(x) - 1 \bigr) .
  \label{eq:13.112}
\end{equation}
Here $g_0$ is the bare value of the coupling constant, while the product of
delta functions, one at each point, enforces the constraint.  Introduce an
integral representation of the delta functions; this requires a second
functional integral over a Lagrange multiplier variable $\alpha(x)$:
\begin{equation}
  \prod_x \delta\bigl( n^2(x) - 1 \bigr)
  = \int \mathcal{D}\alpha\ \exp\Bigl[ \frac{i}{2g_0^2}\int d^dx\
    \alpha(x)\bigl( n^2(x) - 1 \bigr) \Bigr] .
  \label{eq:13.113}
\end{equation}
Now the variable $n$ is unconstrained and appears in the exponent only
quadratically.  Thus, we can integrate over $n$, to obtain
\begin{equation}
  \begin{aligned}
    Z &= \int \mathcal{D}\alpha\ \bigl( \det[-\partial^2 + i\alpha(x)]
      \bigr)^{-N/2}\ \exp\Bigl[ \frac{i}{2g_0^2}\int d^dx\ \alpha(x) \Bigr]\\
    &= \int \mathcal{D}\alpha\ \exp\Bigl[ -\frac{N}{2}\mathrm{tr}\log
      (-\partial^2 + i\alpha) + \frac{i}{2g_0^2}\int d^dx\ \alpha \Bigr] .
  \end{aligned}
  \label{eq:13.114}
\end{equation}
Since we are taking the limit $N\to\infty$ with $g^2N$ held fixed, both terms
in the exponent are of order $N$.  Thus it makes sense to evaluate the
integral by steepest descents.  This entails dominating the integral by the
value of the function $\alpha(x)$ that minimizes the exponent.  To determine
this configuration, we compute the functional derivative of the exponent with
respect to $\alpha(x)$.  This gives the variational equation
\begin{equation}
  \Bigl( \frac{1}{-\partial^2 + i\alpha} \Bigr)(x,x) = \frac{1}{g_0^2N} .
  \label{eq:13.115}
\end{equation}
The left-hand side of this equation must be constant and real; thus, we should
look for a solution in which $\alpha(x)$ is constant and pure imaginary.
Write
\begin{equation}
  i\alpha = m^2 ,
  \label{eq:13.116}
\end{equation}
then $m^2$ obeys
\begin{equation}
  N\ \int \frac{d^dk}{(2\pi)^d}\ \frac{1}{k^2 + m^2}
  = \frac{1}{g_0^2} .
  \label{eq:13.117}
\end{equation}
We will study this equation first in $d = 2$.  If we define the integral in
\eqref{eq:13.117} with a momentum cutoff, we can evaluate this integral and
find the equation for $m$:
\begin{equation}
  \frac{N}{2\pi}\ \log\frac{\Lambda}{m} = \frac{1}{g_0^2} .
  \label{eq:13.118}
\end{equation}
We can make this equation finite by the renormalization
\begin{equation}
  \frac{1}{g_0^2} = \frac{1}{g^2} + \frac{N}{2\pi}\log\frac{\Lambda}{M} ,
  \label{eq:13.119}
\end{equation}
which introduces an arbitrary renormalization scale $M$.  Then we can solve
for $m$, to find
\begin{equation}
  m = M\ \exp\Bigl( -\frac{2\pi}{N g^2} \Bigr) .
  \label{eq:13.120}
\end{equation}
The mass $m$ is the dynamically generated mass of the $\pi$ fields.  It is a
nonperturbative quantity: it vanishes to all orders in the coupling constant,
since $g^2$ appears in the exponent.  The presence of this mass means that the
Goldstone bosons acquire a mass through nonperturbative effects, destroying
the ordered state at long distances.  This is the exact large-$N$ realization
of the Mermin--Wagner theorem.

In $2 < d < 4$, the same analysis gives a nontrivial fixed point and
dynamically generated mass, confirming the picture obtained from the
renormalization group.  For $d > 4$, the equation \eqref{eq:13.117} has only
the trivial solution $m = 0$, and the ordering survives at large distances.

We have now obtained a complete picture of the behavior of the
$O(N)$-symmetric scalar field theory as a function of the dimension.  For
$d > 4$, the theory is described by the free-field fixed point, and the
critical behavior is that of Landau theory.  At $d = 4$, the theory is
renormalizable and the critical behavior is modified by logarithms.  For
$2 < d < 4$, the theory is controlled by the Wilson--Fisher fixed point, and
the critical exponents differ from their Landau values by computable
corrections.  For $d = 2$, the theory is described by the nonlinear sigma
model, which is asymptotically free and exhibits no ordering at long
distances.  This evolution of the behavior of the model as a function of $d$
illustrates the main point of the previous two chapters: The qualitative
behavior of a quantum field theory is determined not by the fundamental
Lagrangian, but rather by the nature of the renormalization group flow and its
fixed points.  These, in turn, depend only on the basic symmetries that are
imposed on the family of Lagrangians that flow into one another.  This
conclusion signals, at the deepest level, the importance of symmetry
principles in determining the fundamental laws of physics.

\problems

\textbf{13.1} \emph{Correction-to-scaling exponent.}  For critical phenomena in
$4-\epsilon$ dimensions, the irrelevant contributions that disappear most
slowly are those associated with the deviation of the coupling constant
$\lambda$ from its fixed-point value.  This gives the most important
nonuniversal correction to the scaling laws derived in Section~13.1.  By
studying the solution of the Callan-Symanzik equation, show that if the bare
value of $\lambda$ differs slightly from $\lambda^{*}$, the Gibbs free energy
receives a correction
\begin{equation}
  G(M,t) \to G(M,t)\cdot\Bigl( 1 + (\lambda-\lambda^{*})\,
    t^{-\omega}\ \cdots \Bigr) .
\end{equation}
This formula defines a new critical exponent $\omega$, called the
\emph{correction-to-scaling exponent}.  Show that
\begin{equation}
  \omega = \nu\,\beta'\bigl( \lambda^{*} \bigr) .
\end{equation}

\textbf{13.2} \emph{The exponent $\eta$.}  By combining the result of
Problem~10.3 with an appropriate renormalization prescription, show that the
leading term in $\gamma(\lambda)$ in $\phi^4$ theory is
\begin{equation}
  \gamma = \frac{\lambda^2}{12(4\pi)^4} .
\end{equation}
Generalize this result to the $O(N)$-symmetric $\phi^4$ theory to derive
Eq.~\eqref{eq:13.47}.  Compute the leading-order ($\epsilon^2$) contribution
to $\eta$.

\textbf{13.3} \emph{The $CP^N$ model.}  The nonlinear sigma model discussed in
the text can be thought of as a quantum theory of fields that are coordinates
on the unit sphere.  A slightly more complicated space of high symmetry is
complex projective space, $CP^N$.  This space can be defined as the space of
$(N+1)$-dimensional complex vectors $(z_1, \ldots, z_{N+1})$ subject to the
condition
\begin{equation}
  \sum_j |z_j|^2 = 1 ,
\end{equation}
with points related by an overall phase rotation identified, that is,
\begin{equation}
  (e^{i\alpha}z_1, \ldots, e^{i\alpha}z_{N+1}) \quad
  \text{identified with} \quad (z_1, \ldots, z_{N+1}) .
\end{equation}
In this problem, we study the two-dimensional quantum field theory whose
fields are coordinates on this space.
\begin{enumerate}[label=(\alph*)]
  \item One way to represent a theory of coordinates on $CP^N$ is to write a
    Lagrangian depending on fields $z_j(x)$, subject to the constraint, which
    also has the local symmetry
    \begin{equation}
      z_j(x) \to e^{i\alpha(x)} z_j(x) ,
    \end{equation}
    independently at each point $x$.  Show that the following Lagrangian has
    this symmetry:
    \begin{equation}
      \mathcal{L} = \frac{1}{g^2}\ \Bigl[ |\partial_\mu z_j|^2
        - |z_j^{*}\partial_\mu z_j|^2 \Bigr] .
    \end{equation}
    To prove the invariance, you will need to use the constraint on the
    $z_j$, and its consequence
    \begin{equation}
      z_j^{*}\partial_\mu z_j = -(\partial_\mu z_j^{*})z_j .
    \end{equation}
    Show that the nonlinear sigma model for the case $N = 3$ can be
    converted to the $CP^N$ model for the case $N = 1$ by the substitution
    \begin{equation}
      n^i = z^{*}_j(\sigma^i)^j_{\ k}\,z^k ,
    \end{equation}
    where $\sigma^i$ are the Pauli sigma matrices.
  \item To write the Lagrangian in a simpler form, introduce a scalar
    Lagrange multiplier $\lambda$ which implements the constraint and also a
    vector Lagrange multiplier $A_\mu$ to express the local symmetry.  More
    specifically, show that the Lagrangian of the $CP^N$ model is obtained
    from the Lagrangian
    \begin{equation}
      \mathcal{L} = \frac{1}{g^2}\ \Bigl[ |D_\mu z_j|^2 - \lambda\bigl(
        |z_j|^2 - 1 \bigr) \Bigr] ,
    \end{equation}
    where $D_\mu = (\partial_\mu + iA_\mu)$, by functionally integrating over
    the fields $\lambda$ and $A_\mu$.
  \item We can solve the $CP^N$ model in the limit $N\to\infty$ by
    integrating over the fields $z_j$.  Show that this integral leads to the
    expression
    \begin{equation}
      Z = \int \mathcal{D}\lambda\,\mathcal{D}A\ \exp\Bigl[ -iN\,
        \mathrm{tr}\log\bigl( -D^2 - \lambda \bigr) + \frac{iN}{g^2}\int d^2x\
          \lambda \Bigr] ,
    \end{equation}
    where we have kept only the leading terms for $N\to\infty$,
    $g^2N$ fixed.  Using methods similar to those we used for the nonlinear
    sigma model, examine the conditions for minimizing the exponent with
    respect to $\lambda$ and $A_\mu$.  Show that these conditions have a
    solution at $A_\mu = 0$ and $\lambda = m^2 > 0$.  Show that, if $g^2$ is
    renormalized at the scale $M$, $m$ can be written as
    \begin{equation}
      m = M\ \exp\Bigl( -\frac{2\pi}{Ng^2} \Bigr) .
    \end{equation}
  \item Now expand the exponent about $A_\mu = 0$.  Show that the first
    nontrivial term in this expansion is proportional to the vacuum
    polarization of massive scalar fields.  Evaluate this expression using
    dimensional regularization, and show that it yields a standard kinetic
    energy term for $A_\mu$.  Thus the strange nonlinear field theory that we
    started with is finally transformed into a theory of $(N+1)$ massive
    scalar fields interacting with a massless photon.
\end{enumerate}
